% sec-03-general-investigational-plan.tex - General Investigational Plan (final IND)
% Phase 1 PDAC IND: AI Generation. IND v1.0; repository v4.3.0.
\section{General Investigational Plan}

This General Investigational Plan describes, at the level of detail required by
21 CFR \S 312.23(a)(3)(iv), the rationale for the investigation, the indication to
be studied, the general approach to evaluating the investigational articles, the
trials to be conducted during the first year, the estimated number of subjects to
be exposed, and the foreseeable drug-related and device-related risks together with
their mitigations. The investigation is a single, integrated, first-in-human study
of two regulated articles administered in combination: perioperative oral
daraxonrasib (RMC-6236), an investigational pan-RAS(ON) multi-selective
small-molecule inhibitor evaluated under the investigational new drug pathway of
21 CFR part 312, and an on-premises large language model (LLM) directed eight-arm
robotic pancreaticoduodenectomy (Whipple) system, a significant-risk Class II
collaborative device evaluated under the investigational device exemption pathway
of 21 CFR part 812. The study carries a single protocol, a single informed consent
bearing the Physical AI opt-out, and a single three-tier safety governance
structure. The plan that follows is intentionally quantitative so that the reviewer
can assess the design, the dose schedule, the operating characteristics, and the
benefit-risk balance from the numbers themselves \cite{phase1guidance,daraxwhipple}.

This plan is written to be self-contained at the level of the General
Investigational Plan while cross-referencing the modules that elaborate each
element: the Investigator Brochure (\S 5) for the integrated nonclinical and
clinical summary of daraxonrasib and the device, the proposed clinical research
module (\S 6) for the full protocol and informed consent, the Chemistry,
Manufacturing and Control information (\S 7), and the Pharmacology and Toxicology
information (\S 8). Where this plan states a numeric specification, that
specification is the binding value carried forward consistently into those modules
and into the protocol, so the reviewer encounters one set of numbers throughout the
submission rather than several. The General Investigational Plan is, by the design
of 21 CFR \S 312.23(a)(3)(iv), a brief description; here it is expanded to the
quantitative depth that a combined drug-device first-in-human investigation in a
lethal indication warrants, because the benefit-risk judgment the reviewer must make
turns on the exact dose schedule, the exact safety envelope, and the exact
operating characteristics rather than on a qualitative summary
\cite{phase1guidance,onpremwhippl}.

\subsection{Rationale}

The scientific rationale rests on a population of profound unmet need and a
mechanistically matched intervention. Pancreatic ductal adenocarcinoma (PDAC)
remains among the most lethal of solid tumors: the five-year overall survival is
below 13 percent, and the disease is the third leading cause of cancer death in the
United States and the fourth in the European Union \cite{pdac060s2030}. The great
majority of PDAC is driven by an activating KRAS mutation at codon 12, most commonly
KRAS G12D, G12V, or G12R, and until recently this oncogenic dependency was
considered pharmacologically intractable. Daraxonrasib (RMC-6236) is an oral
once-daily RAS(ON) multi-selective (pan-RAS) inhibitor that binds the GTP-bound,
active conformation of mutant RAS in a tri-complex with cyclophilin A and thereby
suppresses downstream effector signaling through the MAPK and PI3K pathways. The
agent received FDA Breakthrough Therapy designation in June 2025 for previously
treated metastatic PDAC harboring a KRAS G12 mutation, and the 300 mg once-daily
dose is the established recommended Phase 2 dose of the RASolute 302 pan-KRAS
program \cite{daraxwhipple,oncotrialllm}. Pairing a genotype-matched systemic agent
with a margin-negative curative-intent resection in the resectable and
borderline-resectable setting is a scientifically coherent strategy for a disease in
which the resection margin is the strongest modifiable determinant of survival.

The biological logic of the pairing deserves to be stated precisely because it
shapes the entire dose schedule. In the resectable and borderline-resectable
setting the principal lever on survival is a microscopically margin-negative (R0)
resection, and the principal threat to that lever is residual disease at the
retroperitoneal and superior-mesenteric-vessel margins. A pan-RAS inhibitor that
suppresses the dominant oncogenic driver in the perioperative window has the
potential to reduce the viable tumor burden at those margins and to address
micrometastatic disease that an operation alone cannot reach, while the operation
removes the bulk disease that systemic therapy alone cannot durably control. The
two modalities are therefore complementary rather than redundant, and neither is a
substitute for the other. The investigational article on the device side is not a
second drug but the instrument that delivers the operation to a higher and more
reproducible standard of margin control and vascular safety, which is why the IND
and the IDE are evaluated as one integrated investigation rather than as two
parallel studies \cite{daraxwhipple,onpremwhippl}.

The second component of the rationale is the surgical platform itself. The Whipple
operation is the most technically demanding curative-intent abdominal procedure, and
its early morbidity is dominated by pancreatic fistula, intra-abdominal hemorrhage,
sepsis, and failure to rescue. The investigational eight-arm robotic system is
engineered to perform this operation within hard cyber-physical limits under
continuous human oversight: eight arms each with seven degrees of freedom (56
degrees of freedom total), 640 sensor channels (80 per arm) sampled at 100 kHz for
force and 10 kHz for other modalities, a positional accuracy of 0.05 mm RMS at
1200 mm/s, per-arm tip-force caps at or below 3 N and a cumulative cross-arm cap at
or below 18 N, a 10 kHz safety heartbeat (a 64-byte frame) with a 100 microsecond
watchdog, arm-park within 50 microseconds, a cross-arm emergency stop within 3 ms to
0.0 N, and a hardware emergency stop within 500 ms consistent with 21 CFR
\S 312.404 \cite{onpremwhippl,daraxwhipple}. 

A vascular no-fly gate continuously
protects five named vessels (the superior mesenteric vein, portal vein, hepatic
artery, celiac axis, and superior mesenteric artery) with a 3.0 mm advisory
soft-warning zone, a 1.0 mm no-fly zone for the superior mesenteric and portal veins
and a 1.5 mm no-fly zone for the arterial structures, and a 5.0 mm hard-stop zone
that triggers an emergency stop. The on-premises LLM operates strictly as a
second-opinion advisory oracle, hash-pinned to a deterministic seed and a repository
commit, isolated outside the vendor kinematic stack, and prohibited from full
autonomy consistent with 21 CFR \S 312.21(e); every advisory passes a ten-gate
verification-before-generation suite returning ACCEPT, BLOCK, or ESCALATE, aligned
with H.R. 9510, the Verification Before Generation in Physical AI Oncology Trials
Act of 2026 \cite{congress9510,clintrustpai}.

The operation is decomposed into eight prespecified operative phases so that the
device-feasibility endpoint can be assessed task by task rather than as an
all-or-nothing whole. Phases P1 through P4 cover the dissection and specimen
removal; phase P5 is the pancreaticojejunostomy reconstructed as a duct-to-mucosa
anastomosis with a ring tension held between 0.30 N and 0.60 N; phase P6 is the
hepaticojejunostomy reconstructed end-to-side with a working tension between 0.20 N
and 0.50 N; phase P7 is the gastrojejunostomy reconstructed antecolic with a tension
between 0.40 N and 0.80 N; and phase P8 closes the case with imaging, hemostasis,
and drain placement. These per-phase tension envelopes are not decorative: each is a
real-time hard constraint enforced by the same force-capping logic that bounds the
per-arm and cross-arm tip force, so a reconstruction that would exceed its phase
envelope is prevented before it can injure tissue. The phase decomposition is also
the unit at which the LLM advisory operates, because each anastomotic phase presents
a discrete decision point at which a verified second opinion is most useful
\cite{onpremwhippl,daraxwhipple}.

The schedule rationale is distinct from, and complementary to, the clinical
rationale. This IND is itself the highest-value product of a verified single-prompt
authoring method whose purpose is to compress the administrative and preparation
time that separates clinical milestones, without ever altering tumor biology or any
fixed regulatory clock. Figure 8 ranks the six document targets where faster,
validated authoring yields the greatest schedule value. The initial IND and IRB
package, the subject of this build, is the highest-value target because it sits on
the critical path; one input, faster validated authoring, feeds all six targets, and
each contributes to a common compression of administrative and preparation time
measured in months to a year. Crucially, the gains are capped by clinical events and
by fixed review clocks: faster authoring cannot shorten the 30-day FDA review or
manufacture missing toxicology or stability data \cite{paipredict4x,natlpaiplatf}.

\needspace{8\baselineskip}
\begin{mermaidfig}
\node[mmin]    (in) at (0,-4.5)     {Faster, validated\\IND authoring};
\node[mmgoal]  (t1) at (6.0,0)      {1. Initial IND + IRB\\package (this build)};
\node[mmgoal]  (t2) at (6.0,-1.8)   {2. Protocol amendments\\+ synchronized consent};
\node[mmproc]  (t3) at (6.0,-3.6)   {3. Cohort-review package\\after DLT data mature};
\node[mmgoal]  (t4) at (6.0,-5.4)   {4. Complete clinical-hold\\response};
\node[mmaccent](t5) at (6.0,-7.2)   {5. EOP2 briefing +\\Phase 3 protocol};
\node[mmaccent](t6) at (6.0,-9.0)   {6. CSR + NDA/BLA\\modules at lock};
\node[mmgoal]  (out)at (12.0,-4.5)  {Compressed admin /\\prep time\\(months to a year)};
\node[mmdec]   (lim)at (12.0,-8)  {Capped by clinical\\events and fixed\\review clocks};
\draw[mmedge] (in.east) to[out=70,in=180,looseness=0.5]  (t1.west);
\draw[mmedge] (in.east) to[out=45,in=180,looseness=0.5]  (t2.west);
\draw[mmedge] (in.east) to[out=20,in=180,looseness=0.5]  (t3.west);
\draw[mmedge] (in.east) to[out=-20,in=180,looseness=0.5] (t4.west);
\draw[mmedge] (in.east) to[out=-45,in=180,looseness=0.5] (t5.west);
\draw[mmedge] (in.east) to[out=-70,in=180,looseness=0.5] (t6.west);
\draw[mmedge] (t1.east) to[out=0,in=120,looseness=0.6] (out.west);
\draw[mmedge] (t2.east) to[out=0,in=150,looseness=0.6] (out.west);
\draw[mmedge] (t3.east) -- (out.west);
\draw[mmedge] (t4.east) to[out=0,in=200,looseness=0.6] (out.west);
\draw[mmedge] (t5.east) to[out=0,in=225,looseness=0.6] (out.west);
\draw[mmedge] (t6.east) to[out=0,in=250,looseness=0.6] (out.west);
\draw[mmedged] (out.south) to[out=-90,in=90] (lim.north);
\end{mermaidfig}
\vspace{-0.75cm}
\figcaption{Figure 8. The six greatest-acceleration IND document targets, ranked,
with the initial IND and IRB\\package the highest-value target and the subject of
this build. One input, faster validated authoring, feeds\\all six; each contributes
to a common compression of administrative and preparation time of months to\\a year.
The schedule value is highest on the critical path, and the gains are capped by
clinical events\\and fixed review clocks (they cannot shorten the 30-day FDA review
or manufacture missing data).}

The mechanism of the schedule benefit is specific and is shown in Figure 9. A
clinical-trial timeline divides into three buckets: a clinical and operational
bucket (recruitment, surgery, the dose-limiting-toxicity window, and
progression-free and overall-survival maturation), which is not compressible; an
administrative and preparation bucket (data cleaning, analysis, and document
authoring), which is compressible; and a fixed regulatory-review bucket (the 30-day
IND review and the 30-day clinical-hold review), which is fixed. The repository LLM
acts on the middle bucket only. The exact per-document savings are: the initial IND
and IRB package falls from 8 to 12 weeks to 1 to 4 days; a protocol amendment with
synchronized consent falls from 2 to 4 weeks to 1 to 2 days; a cohort-review package
falls from 1 to 2 weeks to about 1 day; a complete clinical-hold response falls from
3 to 6 weeks to 2 to 4 days; an annual development safety update report (DSUR) falls
from 4 to 8 weeks to 2 to 3 days; and a clinical study report under ICH E3 falls
from 3 to 6 months to 1 to 2 weeks, summing to months to a year saved cumulatively
\cite{paipredict4x,paisitedocpk}.

\needspace{8\baselineskip}
\begin{mermaidfig}
\node[mmdec]  (b1) at (0,0)        {Clinical /\\operational NOT\\ compressible};
\node[mmaccent](b2)at (5.6,0)      {Administrative /\\preparation\\COMPRESSIBLE};
\node[mmdec]  (b3) at (11.2,0)     {Regulatory review\\FIXED};
\node[mmctx]  (b1d)at (0,-2.4)     {Recruitment, surgery,\\DLT window, PFS / OS\\maturation};
\node[mmctx]  (b3d)at (11.2,-2.4)  {30-day IND review,\\30-day clinical-hold\\review};
\node[mmproc] (llm)at (5.6,-2.4)   {Repository LLM acts\\on this bucket only};
\node[mmin]   (s1) at (1.4,-5.0)   {IND + IRB package:\\8-12 weeks to 1-4 days};
\node[mmin]   (s2) at (5.6,-5.0)   {Protocol amendment +\\consent: 2-4 weeks to\\1-2 days};
\node[mmin]   (s3) at (9.8,-5.0)   {Cohort-review package:\\1-2 weeks to ~1 day};
\node[mmin]   (s4) at (1.4,-7.4)   {Clinical-hold response:\\3-6 weeks to 2-4 days};
\node[mmin]   (s5) at (5.6,-7.4)   {Annual DSUR:\\4-8 weeks to 2-3 days};
\node[mmin]   (s6) at (9.8,-7.4)   {CSR (ICH E3):\\3-6 months to 1-2 weeks};
\node[mmgoal] (g)  at (5.6,-9.6)   {Months to a year saved\\cumulatively};
\draw[mmedge] (b1) -- (b1d);
\draw[mmedge] (b3) -- (b3d);
\draw[mmedge] (b2) -- (llm);
\draw[mmedge] (llm.south) to[out=-120,in=90,looseness=0.6] (s1.north);
\draw[mmedge] (llm.south) -- (s2.north);
\draw[mmedge] (llm.south) to[out=-60,in=90,looseness=0.6] (s3.north);
\draw[mmedge] (s1) -- (s4);
\draw[mmedge] (s2) -- (s5);
\draw[mmedge] (s3) -- (s6);
\draw[mmedge] (s4.south) to[out=-90,in=150,looseness=0.6] (g.west);
\draw[mmedge] (s5.south) -- (g.north);
\draw[mmedge] (s6.south) to[out=-90,in=30,looseness=0.6] (g.east);
\end{mermaidfig}
\vspace{-0.75cm}
\figcaption{Figure 9. The three timeline buckets. Only the administrative and
preparation bucket compresses;\\the clinical and operational bucket and the fixed
regulatory review clocks do not. The repository LLM\\acts on the middle bucket only,
with exact per-document savings: initial IND and IRB package 8 to 12\\weeks to 1 to 4
days; protocol amendment with synchronized consent 2 to 4 weeks to 1 to 2 days;\\
cohort-review package 1 to 2 weeks to about 1 day; complete clinical-hold response 3
to\\6 weeks to 2 to 4 days; annual DSUR 4 to 8 weeks to 2 to 3 days; clinical study
report\\3 to 6 months to 1 to 2 weeks, summing to months to a year saved
cumulatively.}

To make the bucket argument auditable rather than rhetorical, Table 4.1 tabulates
the same six documents with their traditional and method-assisted durations, the
implied compression factor, and the bucket each occupies. The compression factor is
computed at the midpoint of each stated range so that the reader can reproduce it:
for the initial IND and IRB package the traditional midpoint is 10 weeks (70 days)
and the method-assisted midpoint is about 2.5 days, a factor near 28; for the
clinical study report the traditional midpoint is about 4.5 months (135 days) and
the method-assisted midpoint is about 10.5 days, a factor near 13. Every entry in
the table sits in the administrative and preparation bucket, which is the only
bucket the method touches; the clinical and operational bucket and the fixed
regulatory-review bucket are listed in the final rows precisely to show that they
carry a compression factor of one by construction \cite{paipredict4x,paisitedocpk}.

\needspace{10\baselineskip}
\noindent\begin{tabularx}{\textwidth}{L{4.1cm} L{2.5cm} L{2.4cm} C{1.6cm} Y}
\toprule
\textbf{Document or interval} & \textbf{Traditional} & \textbf{Method-assisted} & \textbf{Factor} & \textbf{Bucket} \\
\midrule
Initial IND + IRB package (this build) & 8 to 12 weeks & 1 to 4 days & $\approx$28x & Administrative / preparation (compressible) \\
\addlinespace
Protocol amendment + synchronized consent & 2 to 4 weeks & 1 to 2 days & $\approx$14x & Administrative / preparation (compressible) \\
\addlinespace
Cohort-review package (post-DLT) & 1 to 2 weeks & about 1 day & $\approx$11x & Administrative / preparation (compressible) \\
\addlinespace
Complete clinical-hold response & 3 to 6 weeks & 2 to 4 days & $\approx$11x & Administrative / preparation (compressible) \\
\addlinespace
Annual DSUR & 4 to 8 weeks & 2 to 3 days & $\approx$17x & Administrative / preparation (compressible) \\
\addlinespace
Clinical study report (ICH E3) & 3 to 6 months & 1 to 2 weeks & $\approx$13x & Administrative / preparation (compressible) \\
\midrule
30-day IND review & 30 days & 30 days & 1x & Regulatory review (fixed) \\
\addlinespace
Recruitment, surgery, DLT window, PFS / OS maturation & not compressible & not compressible & 1x & Clinical / operational (not compressible) \\
\bottomrule
\end{tabularx}
\vspace{-0.4cm}
\figcaption{Table 4.1. Per-document acceleration with the compression factor
computed at the midpoint of each stated range. Every accelerated item sits in the
administrative and preparation bucket; the fixed regulatory-review clock and the
clinical and operational interval carry a compression factor of one by construction
and are shown for contrast.}

These two rationales meet at the patient. Figure 10 traces the patient time-saved
cascade: faster validated authoring advances each step from IND and IRB submission
and site activation, through the elapse of the 30-day clock and the first dose, to
faster 3+3 cohort turnover toward the recommended Phase 2 dose, shortening the white
space between phases and, for this low-survival population, extending
progression-free and overall survival, without changing tumor biology or any fixed
review clock. The cascade is the patient-level expression of the schedule benefit
and is the bridge between the rationale here and the favorable benefit-risk balance
developed in \S 4.6 \cite{paipredict4x,pdac060s2030}.

\needspace{8\baselineskip}
\begin{mermaidfig}
\node[mmaccent](a) at (0,0)     {Faster validated\\IND authoring};
\node[mmproc] (b) at (4.2,0)    {Earlier IND / IRB\\submission +\\activation};
\node[mmproc] (c) at (8.4,0)    {Earlier 30-day clock\\elapses; first dose};
\node[mmproc] (d) at (12.6,0)   {Faster 3+3 cohort\\turnover to RP2D};
\node[mmproc] (e) at (12.6,-2.3){Shorter white space\\between phases};
\node[mmgoal] (f) at (8.4,-2.3) {Earlier access to\\effective therapy};
\node[mmgoal] (g) at (4.2,-2.3) {Extended PFS / OS\\(5-year OS < 13\%)};
\node[mmctx]  (n) at (8.4,-4.2) {Does not change tumor biology\\or fixed review clocks};
\draw[mmedge] (a) -- (b);
\draw[mmedge] (b) -- (c);
\draw[mmedge] (c) -- (d);
\draw[mmedge] (d) -- (e);
\draw[mmedge] (e) -- (f);
\draw[mmedge] (f) -- (g);
\draw[mmedged] (n.west) to[out=180,in=-90,looseness=0.7] (g.south);
\end{mermaidfig}
\vspace{-0.75cm}
\figcaption{Figure 10. The patient time-saved cascade. Faster validated authoring
advances each step from IND submission to first dose, then through faster 3+3 cohort
turnover to the recommended Phase 2 dose, shortening the white space between phases
and, for this low-survival population (five-year overall survival under 13 percent),
extending progression-free and overall survival, without changing tumor biology or
the fixed regulatory review clocks.}

A worked schedule example fixes the magnitude of the cascade in days. Suppose the
administrative authoring of the initial IND and IRB package is compressed from a
traditional 70 days (the 10-week midpoint) to about 3 days. The fixed 30-day FDA
review clock is unchanged, the operation and the 28-day DLT window are unchanged,
and the survival follow-up clocks are unchanged. The net effect is that the first
dose and the first operation are reached approximately 67 days earlier than they
would otherwise have been, and that 67-day advance propagates through every
downstream cohort because each cohort's start depends on the prior cohort's
completion. Across three cohorts and the inter-phase authoring that connects this
study to a later-phase program, the accumulated advance reaches the months-to-a-year
figure stated in Figure 9. None of this advance is borrowed from biology or from
patient safety: the DLT window, the sentinel staggering, and the review clocks
remain exactly as specified \cite{paipredict4x,clintrustpai}.

\subsection{Indication to be Studied}

The indication to be studied is histologically or cytologically confirmed pancreatic
ductal adenocarcinoma that is resectable or borderline-resectable on
multidisciplinary radiologic review, harbors a documented KRAS G12 mutation (for
example G12D, G12V, or G12R) on a validated tumor genotyping assay, occurs in an
adult of at least 18 years with an Eastern Cooperative Oncology Group (ECOG)
performance status of 0 or 1, and is anatomically compatible with the eight-arm
robotic system and with a planned pancreaticoduodenectomy. Eligibility for
perioperative daraxonrasib follows the broad pan-KRAS criteria of the RASolute 302
program \cite{daraxwhipple}. The indication is deliberately narrow at this
first-in-human stage so that the safety and feasibility readouts are interpretable:
distant metastatic disease, vascular involvement precluding a margin-negative
resection, prior pancreatic resection precluding the planned reconstruction,
uncontrolled intercurrent illness, pregnancy or lactation, and known
contraindication to daraxonrasib are exclusionary. Table 4.2 summarizes the
eligibility frame that defines the indication; the full inclusion and exclusion
criteria are set out in the protocol and reproduced in the proposed clinical
research module of this IND.

\needspace{12\baselineskip}
\noindent\begin{tabularx}{\textwidth}{L{3.0cm} Y Y}
\toprule
\textbf{Domain} & \textbf{Inclusion (must meet all)} & \textbf{Exclusion (any one disqualifies)} \\
\midrule
Disease and stage & Histologically or cytologically confirmed PDAC; resectable or borderline-resectable on multidisciplinary review. & Distant metastatic disease on staging; vascular involvement precluding R0 resection or anatomy incompatible with the eight-arm system. \\
\addlinespace
Genotype & Documented KRAS G12 mutation (G12D, G12V, or G12R) on a validated assay; eligible under pan-KRAS RASolute 302 criteria. & Absence of a qualifying KRAS G12 mutation. \\
\addlinespace
Performance and organ function & ECOG performance status 0 or 1; adequate hematologic, hepatic, and renal function for major pancreaticobiliary surgery and daraxonrasib. & Uncontrolled cardiac, pulmonary, hepatic, renal, or active infectious disease making surgery or dosing unsafe. \\
\addlinespace
Surgical candidacy & Candidate for pancreaticoduodenectomy; tumor and vascular anatomy compatible with port placement and the working volume. & Prior pancreatic resection or prior surgery precluding the planned reconstruction. \\
\addlinespace
Age and consent & Age at least 18 years; signed informed consent including explicit election or declination of the Physical AI opt-out. & Inability to consent or to comply with the Schedule of Activities. \\
\addlinespace
Other & Measured baseline CA 19-9 recorded for serial assessment (no threshold imposed). & Pregnancy or lactation; known contraindication or hypersensitivity to daraxonrasib or a required excipient; unmanageable drug interaction. \\
\bottomrule
\end{tabularx}
\vspace*{-0.25cm}
\figcaption{Table 4.2. Eligibility summary defining the indication to be studied.
Selection is equitable\\and is determined solely by the scientific and safety
criteria; limited English proficiency is\\not an exclusion, and qualified
interpreters and translated consent materials are provided.}

The indication is positioned within the larger PDAC treatment landscape so that the
reviewer can see what is and is not claimed. The resectable and borderline-resectable
population is a minority of all PDAC at presentation, because most patients present
with locally advanced or metastatic disease; it is, however, the population in which
a curative-intent resection is anatomically possible and therefore the population in
which margin control and perioperative systemic therapy can change the survival
trajectory. The genotype restriction to KRAS G12 is mechanistic rather than
arbitrary: daraxonrasib acts on the GTP-bound active conformation of mutant RAS, so
a documented qualifying mutation is a prerequisite for the drug to have its intended
target. The anatomic restriction to tumors compatible with the eight-arm working
volume and with the five-vessel no-fly geometry is a device-safety prerequisite, not
a comment on resectability in general. The indication is thus the intersection of an
oncologic window, a genotype, and an anatomic envelope, and each restriction maps to
a element of intervention \cite{daraxwhipple,onpremwhippl}.

Equity in selection is an explicit design commitment. The eligibility frame imposes
no criterion that is not scientifically or safety-justified; limited English
proficiency is not an exclusion, and qualified interpreters and translated consent
materials are provided so that language is never the reason a patient cannot enroll.
No threshold is imposed on the baseline CA 19-9; it is measured and recorded only so
that the serial trend can be followed, because a substantial fraction of PDAC is CA
19-9 non-secretory and a threshold would exclude those patients without scientific
justification. The screen-to-treat ratio, approximately 36 screened to reach up to
18 treated, reflects the genotype and anatomic restrictions rather than any
socioeconomic filter, and the screening log is retained so that the equity of the
selection can be audited \cite{daraxwhipple,clintrustpai}.

A synthetic reference exemplar, designated PAT-PDAC-0001, illustrates a
representative eligible participant and is used for design illustration only: a
68-year-old with a 3.4 cm head-of-pancreas tumor abutting the superior mesenteric
vein at 75 degrees, KRAS G12D, microsatellite stable, a baseline CA 19-9 of 412
U/mL, ECOG performance status 1, after a completed neoadjuvant modified FOLFIRINOX
course of four cycles. This exemplar does not represent an enrolled individual but
fixes the anatomic and oncologic context in which the eligibility frame and the
device no-fly gating are exercised. The 75-degree venous abutment, for instance, is
exactly the geometry in which the 1.0 mm no-fly zone for the superior mesenteric and
portal veins and the 3.0 mm soft-warning zone are most relevant, because the
dissection plane runs along the vessel for a substantial arc; the exemplar therefore
serves as a concrete test case for the device-safety argument rather than as a
generic patient sketch \cite{daraxwhipple,onpremwhippl}.

\subsection{General Approach for Evaluation of Treatment}

The general approach is a prospective, open-label, single-arm, first-in-human study
that combines a standard Phase 1 3+3 dose-finding evaluation of perioperative
daraxonrasib with an early-feasibility evaluation of the LLM-directed eight-arm
robotic Whipple system. A single-arm, open-label design is scientifically and
ethically appropriate, and a randomized comparator arm is not, because the
first-in-human objectives, characterizing early safety, identifying a recommended
Phase 2 dose, and demonstrating that a novel surgical system can perform its
intended function within hard cyber-physical limits, are not comparative-efficacy
questions and are answered most safely in a small, intensively monitored cohort with
sentinel staggered enrollment and case-by-case data and safety monitoring board
(DSMB) review. The early-feasibility pathway recognized by FDA for significant-risk
devices is explicitly structured around small first-in-human cohorts that generate
the device-performance and safety evidence needed before any controlled comparison
\cite{phase1guidance}. The 2025 Dutch nationwide cohort of 1000 robotic
pancreaticoduodenectomies provides a robust external human benchmark for the
secondary surgical-quality endpoints, so the absence of an internal control does not
leave the surgical results uninterpretable \cite{daraxwhipple}.

The choice of an external benchmark over an internal control is a deliberate design
decision with explicit limits. A nationwide cohort of 1000 consecutive robotic
pancreaticoduodenectomies captures the contemporary standard of robotic Whipple
practice across many surgeons and the learning-curve gradient that accompanies any
new surgical platform, and it supplies stable comparator rates for the four
secondary surgical-quality endpoints that the present study cannot internally
control: an ISGPS grade B/C fistula rate of 24.4 percent, a Clavien-Dindo grade III
or higher complication rate of 41.3 percent, an unplanned conversion rate of 10.1
percent, and a 30-day mortality of 3.9 percent. These comparators are used for
descriptive contextualization only. The single-arm design supports estimation of the
study rates with exact confidence intervals; it does not support inferential
between-group testing against the Dutch cohort, because the populations, the
case mix, and the perioperative systemic therapy differ, and because a first-in-human
cohort of up to 18 is far too small to power such a comparison. The benchmark
therefore answers the question "is the observed rate in the expected neighborhood of
contemporary practice" and not the question "is this system superior"
\cite{daraxwhipple,phase1guidance}.

\subsubsection{Objectives and Endpoints}

The study carries three co-primary objectives, one for each axis of the combined
IND/IDE investigation, beneath which sit secondary oncologic-surgical quality
endpoints benchmarked against the Dutch 2025 cohort and exploratory non-confirmatory
endpoints. The first co-primary objective characterizes the early safety of the
combined intervention; its endpoint is the incidence of device-related or
procedure-related serious adverse events through 30 days after the procedure,
including Clavien-Dindo grade III or higher complications, with device or procedure
attribution adjudicated by an independent reviewer blinded to the daraxonrasib dose
level. 

The second co-primary objective determines the maximum tolerated dose (MTD)
and recommended Phase 2 dose (RP2D) of perioperative daraxonrasib, defined
operationally by the rate of dose-limiting toxicities (DLTs) within the 28-day DLT
window of the 3+3 escalation. The third co-primary objective establishes the
technical feasibility of the eight-arm robotic Whipple under continuous human
oversight; its endpoint is the proportion of procedures that complete all
prespecified assigned operative tasks without unplanned conversion caused by device
malfunction or unsafe device behavior. Figure 11 renders the objective-to-endpoint
hierarchy, and Table 4.3 sets out the full mapping with justifications and the Dutch
2025 comparators \cite{daraxwhipple,phase1guidance}.

\needspace{8\baselineskip}
\begin{mermaidfig}
\node[mmgoal] (pr) at (0,0)       {PRIMARY (co-primary)\\safety + feasibility};
\node[mmstep] (p1) at (-4.6,-2.2) {Device / procedure SAE\\incidence through 30 days};
\node[mmstep] (p2) at (0,-2.2)    {MTD / RP2D of\\daraxonrasib\\(3+3, 28-day DLT)};
\node[mmstep] (p3) at (4.6,-2.2)  {Assigned tasks completed\\without unsafe\\conversion};
\node[mmaccent](se)at (0,-4.8)    {SECONDARY\\oncologic-surgical};
\node[mmin]   (s1) at (-4.6,-6.45) {R0 rate; ISGPS\\B/C fistula};
\node[mmin]   (s2) at (0,-6.45)    {Clavien-Dindo III+;\\30 / 90-day mortality};
\node[mmin]   (s3) at (4.6,-7)  {Conversion rate; operative\\parameters (vs Dutch 2025)};
\node[mmctx]  (ex) at (0,-8.2)   {EXPLORATORY\\non-confirmatory};
\node[mmin]   (e1) at (0,-10.2)   {PFS / OS to 24 mo (RECIST);\\LLM concordance;\\sim-to-real gap};
\draw[mmedge] (pr) -- (p1);
\draw[mmedge] (pr) -- (p2);
\draw[mmedge] (pr) -- (p3);
\draw[mmedged] (pr) -- (se);
\draw[mmedge] (se) -- (s1);
\draw[mmedge] (se) -- (s2);
\draw[mmedge] (se) -- (s3);
\draw[mmedged] (se.west) to[out=180,in=180,looseness=0.6] (ex.west);
\draw[mmedge] (ex) -- (e1);
\end{mermaidfig}
\vspace*{-0.7cm}
\figcaption{Figure 11. Three co-primary objectives (device or procedure
serious-adverse-event incidence through 30 days, the maximum tolerated dose and
recommended Phase 2 dose of daraxonrasib by 3+3 over the 28-day window, and
completion of assigned operative tasks without unsafe conversion) sit above
secondary oncologic-surgical quality endpoints (R0 rate, ISGPS grade B/C fistula,
Clavien-Dindo grade III or higher, 30 and 90-day mortality, conversion rate, and
operative parameters benchmarked against the Dutch 2025 robotic cohort).}

\needspace{12\baselineskip}
\noindent\begin{tabularx}{\textwidth}{L{2.4cm} L{6.7cm} Y}
\toprule
\textbf{Tier} & \textbf{Objective and endpoint} & \textbf{Justification / comparator} \\
\midrule
\multicolumn{3}{l}{\textit{Co-primary}} \\
\midrule
Safety & Device-related or procedure-related SAE incidence through 30 days, including Clavien-Dindo grade III+, dose-blinded attribution. & Standard first-in-human safety anchor; captures the early lethal cascade of fistula, hemorrhage, sepsis, and failure to rescue. \\
\addlinespace
Dose-finding & MTD and RP2D of perioperative daraxonrasib by the DLT rate within 28-day window of 3+3 across DL1 to DL3. & Established dose-finding design: targeted oncology agents; anchored to RASolute 302 pan-KRAS. \\
\addlinespace
Feasibility & Procedures completing assigned operative tasks, no unplanned conversion from device malfunction or unsafe behavior. & FDA-recognized early-feasibility readout for a significant-risk device. \\
\midrule
\multicolumn{3}{l}{\textit{Secondary (descriptive, Dutch 2025 comparators)}} \\
\midrule
Oncologic adequacy & R0 (microscopically margin-negative) resection rate at day-90 pathology. & Strongest modifiable surgical determinant of survival; descriptive comparison with Dutch 2025. \\
\addlinespace
Fistula & ISGPS grade B/C postoperative pancreatic fistula rate. & Leading driver of post-Whipple morbidity; Dutch 2025 comparator 24.4 percent. \\
\addlinespace
Major complications & Clavien-Dindo grade III+ complication rate through 90 days. & Validated severity threshold; Dutch 2025 comparator 41.3 percent. \\
\addlinespace
Mortality & 30-day and 90-day all-cause mortality. & Short-term survival safety; Dutch 2025 30-day comparator 3.9 percent. \\
\addlinespace
Conversion & Unplanned conversion rate from robotic system to open or conventional approach. & Contextualizes feasibility; Dutch 2025 comparator 10.1 percent. \\
\addlinespace
Operative parameters & Estimated blood loss, operative time, and length of stay. & Mean and median with range; Dutch 2025 learning-curve context. \\
\midrule
\multicolumn{3}{l}{\textit{Exploratory (non-confirmatory)}} \\
\midrule
Survival & PFS and OS assessed to 24 months by RECIST v1.1, Kaplan-Meier. & Preliminary, hypothesis-generating signal; the small single-arm sample precludes confirmatory inference. \\
\addlinespace
Advisory reliability & Concordance between the LLM advisory output and the independent hard-coded safety gate. & Characterizes second-opinion oracle, informs autonomy-graduation decisions reserved for later phases. \\
\addlinespace
Simulation fidelity & Sim-to-real gap in trajectory (mm) and tip force (N). & Phase 0 simulation validation gate; targets < 2 mm, < 0.5 N. \\
\bottomrule
\end{tabularx}
\vspace*{-0.4cm}
\figcaption{Table 4.3. The objective-to-endpoint mapping with justifications and the
Dutch 2025\\nationwide robotic-pancreatoduodenectomy comparators. Comparisons are
descriptive only.}

The three co-primary objectives are weighted equally rather than ranked, because the
investigation cannot be declared a success on any single axis alone. A favorable
drug-safety and dose-finding result is uninformative if the device cannot complete
the operation within its safety envelope, and an excellent device-feasibility result
is unacceptable if the perioperative drug schedule is intolerable. The co-primary
structure therefore requires that the early-safety, dose-finding, and feasibility
readouts each be satisfactory before the combined intervention is judged ready to
advance, and it is the reason the safety-attribution adjudication is performed
blinded to the dose level, so that a device or procedure event is not silently
attributed to the drug or the reverse. The secondary endpoints are descriptive and
support the interpretation of the co-primary results; the exploratory endpoints are
explicitly hypothesis-generating and are reported with the caveat that the sample
size precludes confirmatory inference \cite{phase1guidance,ichpaistduni}.

\subsubsection{Dose-Escalation Approach}

Daraxonrasib is administered orally once daily. The escalation uses exact doses
anchored to the RASolute 302 pan-KRAS program, in which the 300 mg once-daily dose
is the established RP2D. To preserve a conservative first-in-human margin in the
perioperative setting, the escalation begins below that reference dose and steps up
to it: DL1 is 160 mg once daily (the starting dose), DL2 is 220 mg once daily, and
DL3 is the 300 mg once-daily reference dose. Three participants enroll at each level
and are observed over a 28-day DLT window. If no DLT occurs among the first three,
the cohort escalates; if one of three experiences a DLT, the cohort expands to six;
a single DLT among six permits escalation, whereas two or more DLTs at a level
define that level as exceeding the MTD and fix the MTD at the preceding level. The
MTD is the highest level at which fewer than two of six DLT-evaluable participants
experience a DLT, and the RP2D is selected at or below the MTD with integration of
all safety, tolerability, and available pharmacokinetic data. Figure 12 renders the
decision automaton and Table 4.4 states the rule explicitly for each cohort outcome
\cite{daraxwhipple,phase1guidance}.

\needspace{8\baselineskip}
\begin{mermaidfig}
\node[mmin]  (dl1) at (0,0)       {DL1: 160 mg QD\\enroll 3};
\node[mmdec] (g1)  at (0,-2.6)    {DLT in\\28-day window?};
\node[mmstep](e1)  at (4.6,-2.6)  {Expand to 6\\at this level};
\node[mmin]  (dl2) at (0,-5.4)    {DL2: 220 mg QD\\enroll 3};
\node[mmdec] (g2)  at (0,-8.0)    {DLT?};
\node[mmstep](e2)  at (4.6,-8.0)  {Expand to 6};
\node[mmin]  (dl3) at (0,-10.8)   {DL3: 300 mg QD\\enroll 3};
\node[mmdec] (g3)  at (0,-13.4)   {DLT?};
\node[mmstep](e3)  at (4.6,-13.4) {Expand to 6};
\node[mmdark](mtd) at (9.4,-8.0)  {MTD = prior level\\($\geq$ 2 of 6 DLT)};
\node[mmgoal](rp)  at (9.4,-13.35) {RP2D at or below MTD\\integrating PK + safety};
\draw[mmedge] (dl1) -- (g1);
\draw[mmedge] (g1) -- node[left,font=\scriptsize]{0 of 3} (dl2);
\draw[mmedge] (g1) -- node[above,font=\scriptsize]{1 of 3} (e1);
\draw[mmedge] (e1) to[out=-90,in=0,looseness=0.6] node[right,font=\scriptsize,pos=0.2]{1 of 6} (dl2.east);
\draw[mmedge] (e1) to[out=0,in=90,looseness=0.7] node[right,font=\scriptsize,pos=0.3]{$\geq$ 2 of 6} (mtd.north);
\draw[mmedge] (dl2) -- (g2);
\draw[mmedge] (g2) -- node[left,font=\scriptsize]{0 of 3} (dl3);
\draw[mmedge] (g2) -- node[above,font=\scriptsize]{1 of 3} (e2);
\draw[mmedge] (e2) to[out=-90,in=50,looseness=0.6] node[left,font=\scriptsize,pos=0.3]{1 of 6} (dl3.east);
\draw[mmedge] (e2) -- node[above,font=\scriptsize]{$\geq$ 2 of 6} (mtd);
\draw[mmedge] (dl3) -- (g3);
\draw[mmedge] (g3.south) -- node[above,font=\scriptsize]{0 of 3 / 1 of 6} (rp.south west);
\draw[mmedge] (e3) to[out=90,in=-90,looseness=0.7] node[right,font=\scriptsize,pos=0.4]{$\geq$ 2 of 6} (mtd.south);
\draw[mmedge] (g3) -- node[above,font=\scriptsize,pos=0.6]{1 of 3} (e3);
\draw[mmedge] (mtd) -- (rp);
\end{mermaidfig}
\vspace*{-0.7cm}
\figcaption{Figure 12. The standard 3+3 dose-escalation automaton across the three
daraxonrasib dose levels (DL1 160 mg, DL2 220 mg, DL3 300 mg once daily). Three
participants enroll at each level; with no dose-limiting toxicity in the 28-day
window the cohort escalates, with one of three it expands to six, and two or more of
six fixes the maximum tolerated dose at the prior level. The recommended Phase 2 dose
is selected at or below the maximum tolerated dose, integrating pharmacokinetic and
safety data. The design maximum is three levels times six, up to 18 treated.}

\needspace{12\baselineskip}
\noindent\begin{tabularx}{\textwidth}{L{2.6cm} L{2.4cm} L{3.4cm} Y}
\toprule
\textbf{Cohort observation} & \textbf{DLT count} & \textbf{Decision} & \textbf{Consequence} \\
\midrule
First 3 at a level & 0 of 3 & Escalate & Open next higher dose level; current level deemed tolerated at cohort of 3. \\
\addlinespace
First 3 at a level & 1 of 3 & Expand to 6 & Enroll 3 additional participants at same level before any escalation decision. \\
\addlinespace
First 3 at a level & $\geq$ 2 of 3 & MTD exceeded & Do not enroll further at this level; the MTD is the preceding (lower) level. \\
\addlinespace
Expanded 6 at a level & 1 of 6 & Escalate & Single DLT in six permits escalation to the next higher level. \\
\addlinespace
Expanded 6 at a level & $\geq$ 2 of 6 & MTD exceeded & The level exceeds the MTD; the MTD is the preceding level. \\
\addlinespace
Highest level (DL3) & 0 of 3 or 1 of 6 & RP2D candidate & DL3 (300 mg) tolerated; RP2D selected at or below DL3 integrating PK, safety. \\
\addlinespace
Any level & $\geq$ 2 of 6 & Declare MTD / RP2D & MTD fixed at prior level; RP2D at or below MTD. \\
\bottomrule
\end{tabularx}
\vspace*{-0.3cm}
\figcaption{Table 4.4. The 3+3 dose-escalation decision rule. The maximum tolerated
dose is the highest level at which fewer than two of six DLT-evaluable participants
experience a dose-limiting toxicity within the 28-day window; the recommended Phase
2 dose is selected at or below the MTD. Design maximum is three levels times six,
up to 18.}

The dose spacing is itself a safety choice. The step from DL1 to DL2 is 60 mg (a
1.375-fold increment) and the step from DL2 to DL3 is 80 mg (a 1.364-fold
increment), so neither escalation exceeds roughly a 1.4-fold increase in exposure,
which is conservative for a targeted agent whose RP2D is already established. Because
DL3 is the established RASolute 302 RP2D rather than an extrapolated supratherapeutic
dose, the escalation never asks a participant to receive a dose above the dose at
which the agent's tolerability is already characterized in the metastatic setting;
the perioperative novelty is the timing and the surgical context, not an
unprecedented exposure. This is the reason the starting dose is set below the
reference RP2D rather than at it: in the perioperative setting the additional
hazard of impaired wound healing and anastomotic integrity justifies beginning with
a conservative margin and confirming tolerability before stepping up to the
reference dose \cite{daraxwhipple,oncotrialllm}.

A dose-limiting toxicity is defined for the 28-day window as a clinically significant
toxicity attributable to daraxonrasib that meets the protocol's prespecified grade
and persistence thresholds, including a perioperative event in which the drug is the
probable contributor to impaired anastomotic healing. The 28-day window is chosen to
span the early postoperative healing interval during which a drug-related effect on
the pancreaticojejunal and hepaticojejunal anastomoses would manifest, and it
aligns the drug-toxicity assessment with the same 30-day horizon over which the
device-and-procedure serious-adverse-event endpoint is captured. The DLT-evaluable
population comprises participants who either complete the full 28-day window or
experience a DLT within it; a participant who is not evaluable for a non-toxicity
reason is replaced so that each dose level reaches its required number of
DLT-evaluable participants, which preserves the operating characteristics of the
3+3 rule \cite{phase1guidance,ichpaistduni}.

The perioperative pause-and-restart of daraxonrasib is an integral part of the dose
schedule rather than a separate procedure. The drug is held before surgery, and the
restart day is set by an LLM-bound, human-confirmed advisory whose inputs are the
postoperative pancreatic fistula grade by the International Study Group of Pancreatic
Surgery (ISGPS) classification (grade A, B, or C) and the serum daraxonrasib trough
relative to a 0.5 ng/mL threshold, with the four permitted outputs T+7 days, T+14
days, T+21 days, or hold. A 32-iteration deterministic sweep of this advisory
returned a baseline serum of 0.45 ng/mL (sub-threshold) in all iterations, a T+7-day
restart in 29 of 32 iterations (90.6 percent, corresponding to grade A and
sub-threshold), a T+14-day restart in 3 of 32 iterations (9.4 percent, corresponding
to grade B or a delayed trough), and no T+21-day restart. The sweep is reproduced in
Table 4.5 so that the reviewer can see that the advisory's restart-day distribution
is dominated by the earliest safe restart and that no iteration produced a
prolonged hold, while the human confirmation step preserves clinician authority over
every restart decision \cite{daraxwhipple,clintrustpai}.

\needspace{10\baselineskip}
\noindent\begin{tabularx}{\textwidth}{L{3.2cm} C{2.0cm} C{2.4cm} Y}
\toprule
\textbf{Advisory output} & \textbf{Iterations} & \textbf{Share} & \textbf{Driving condition} \\
\midrule
Restart at T+7 days & 29 of 32 & 90.6\% & ISGPS grade A fistula and serum trough sub-threshold (0.45 ng/mL $<$ 0.5 ng/mL). \\
\addlinespace
Restart at T+14 days & 3 of 32 & 9.4\% & ISGPS grade B fistula or a delayed serum trough; restart deferred one week. \\
\addlinespace
Restart at T+21 days & 0 of 32 & 0.0\% & No iteration met the grade or trough condition for a three-week deferral. \\
\addlinespace
Hold (no restart) & 0 of 32 & 0.0\% & No iteration met the hold condition; every iteration restarted within three weeks. \\
\bottomrule
\end{tabularx}
\vspace*{-0.35cm}
\figcaption{Table 4.5. The 32-iteration deterministic sweep of the perioperative
daraxonrasib pause-and-restart advisory. The baseline serum trough was 0.45 ng/mL
(sub-threshold) in every iteration. The restart-day distribution is dominated by the
earliest safe restart (T+7 days in 90.6 percent), and every advisory is
human-confirmed before acted upon.}

\subsubsection{Device Readiness and Statistical Framework}

The device has no pharmacologic dose; its dose analogue is the readiness evidence
required before any patient-facing robotic activity, escalated only by passing a
mandatory Phase 0 simulation-validation gate. Three conditions must be met and
documented before the first robotic case: a Unified Safety Level (USL) rating of at
least 7.0 on the surgical-readiness scale; completion of at least 1000 simulated
procedures across at least two independent frameworks (for example Isaac, MuJoCo,
Gazebo, or PyBullet), so the finding is not an artifact of any single simulator; and
a sim-to-real fidelity within prespecified limits, a trajectory gap less than 2 mm
and a tip-force gap less than 0.5 N between the simulation prediction and the
operative record. The gate is re-validated after any change to the system. These
limits operate together with the intra-operative hard caps (a cross-arm
emergency-stop budget at or below 3 ms, per-arm tip-force caps at or below 3 N, a
cumulative cross-arm cap at or below 18 N, and vascular no-fly gating) so the device
dose analogue is met only when both the simulation evidence and the real-time safety
envelope are demonstrably in place \cite{onpremwhippl,daraxwhipple}.

The cyber-physical safety envelope is specified numerically so that each limit is
verifiable rather than aspirational, and Table 4.6 collects the binding values in one
place. The temporal hierarchy of the safety responses is itself meaningful: the
10 kHz heartbeat samples the system every 100 microseconds and a 100 microsecond
watchdog detects a missed frame within one cycle; an arm parks within 50
microseconds; a cross-arm emergency stop drives all arms to 0.0 N within 3 ms; and a
hardware emergency stop completes within 500 ms. Each successive response operates at
a coarser time scale and a higher level of authority, so a fault is caught at the
finest scale at which it can be caught and escalates only if it is not resolved.
The vascular no-fly gate adds a spatial hierarchy on top of this temporal one: a
3.0 mm soft-warning zone issues an advisory, a 1.0 mm (veins) or 1.5 mm (arteries)
no-fly zone blocks motion toward the vessel, and a 5.0 mm hard-stop zone triggers an
emergency stop. In a representative 100-tick sampling window the gate returned 83
clear verdicts, 6 soft-warning verdicts, 6 no-fly verdicts, and 5 hard-stop verdicts,
which illustrates that the gate is active and discriminating rather than dormant
\cite{onpremwhippl,daraxwhipple}.

\needspace{12\baselineskip}
\noindent\begin{tabularx}{\textwidth}{L{4.4cm} L{5.6cm} Y}
\toprule
\textbf{Parameter} & \textbf{Binding value} & \textbf{Function} \\
\midrule
Degrees of freedom & 8 arms x 7 = 56 & Working envelope for the eight-arm collaborative system. \\
\addlinespace
Sensor channels & 640 (80 per arm) & 100 kHz force sampling; 10 kHz for other modalities. \\
\addlinespace
Positional accuracy & 0.05 mm RMS at 1200 mm/s & Trajectory precision under motion. \\
\addlinespace
Per-arm tip-force cap & $\leq$ 3 N & Bounds the force any single arm applies to tissue. \\
\addlinespace
Cumulative cross-arm cap & $\leq$ 18 N & Bounds the total force applied across all arms. \\
\addlinespace
Safety heartbeat / watchdog & 10 kHz (64-byte frame) / 100 microsecond watchdog & Detects a missed frame within one cycle. \\
\addlinespace
Arm-park latency & $\leq$ 50 microseconds & Parks an arm on a local fault. \\
\addlinespace
Cross-arm emergency stop & $\leq$ 3 ms to 0.0 N & Drives all arms to zero force. \\
\addlinespace
Hardware emergency stop & $\leq$ 500 ms (21 CFR \S 312.404) & Final hardware-level stop authority. \\
\addlinespace
Vascular no-fly zones & 3.0 mm warn; 1.0 mm vein / 1.5 mm artery no-fly; 5.0 mm hard stop & Protects SMV, PV, HA, CA, SMA. \\
\bottomrule
\end{tabularx}
\vspace*{-0.35cm}
\figcaption{Table 4.6. The binding cyber-physical safety envelope of the eight-arm
robotic Whipple system. Each value is a hard, verifiable limit; the temporal
responses form a hierarchy from a 100 microsecond watchdog to a 500 ms hardware
emergency stop, and the vascular no-fly gate adds a spatial hierarchy of warn,
no-fly, and hard-stop zones.}

The statistical framework is that of a small, intensively monitored, single-arm
investigation whose primary purpose is to characterize early safety and to identify
a recommended dose, not to test a comparative-efficacy hypothesis. The analyses are
predominantly descriptive and estimation-based, follow ICH E9, and are reported in
accordance with TRIPOD+AI for the model-assisted advisory component
\cite{ichpaistduni}. No power calculation is performed because the study tests no
confirmatory efficacy hypothesis. Every proportion is accompanied by an exact
(Clopper-Pearson) two-sided 95 percent confidence interval rather than a
normal-approximation interval, which is unreliable near the bounds and at small
sample sizes; time-to-event variables are summarized by Kaplan-Meier with the number
at risk; missing data are not imputed for the primary or safety analyses; and a
nominal alpha of 0.05 two-sided applies to supportive estimation only, with no
multiplicity adjustment because the analyses are descriptive and non-confirmatory.
Four analysis populations are applied consistently: the Safety population (all who
received any intervention), the DLT-evaluable population (all dosed participants who
completed the 28-day window or had a DLT within it), the Per-Protocol population,
and the modified intention-to-treat population \cite{ichpaistduni,phase1guidance}.

The choice of the exact (Clopper-Pearson) interval over the normal-approximation
(Wald) interval is consequential at these sample sizes and is therefore stated
explicitly. The Wald interval is constructed as a symmetric band around the point
estimate and, when the observed proportion is zero or one, collapses to a degenerate
interval of zero width that falsely implies certainty; at a denominator of three or
six this failure is not a rare edge case but the expected situation, because a clean
safety record is the hoped-for result. The Clopper-Pearson interval, by inverting the
exact binomial tail probabilities, always returns a non-degenerate interval that
honestly reflects how little a small denominator can exclude. The TRIPOD+AI reporting
discipline is applied to the LLM advisory component so that the advisory's inputs,
outputs, deterministic seed, pinned commit, and concordance with the hard-coded gate
are all reported in a form that an external reviewer can audit, which is the
model-facing analogue of the participant-level data tabulation applied to the
clinical results \cite{ichpaistduni,clintrustpai}.

\subsection{Description of First Year Trial(s)}

During the first year, a single trial is conducted at a single academic medical
center: the combined IND/IDE first-in-human study of perioperative daraxonrasib and
the LLM-directed eight-arm robotic Whipple. Enrollment is planned for up to 18
treated participants, allocated across the three daraxonrasib dose levels (DL1
through DL3) by the 3+3 rule with staggered sentinel enrollment, from approximately
36 screened. Recruitment proceeds at approximately one to two participants per month
through the institutional multidisciplinary pancreatic-cancer program, with eligible
individuals identified at tumor-board review and described to them by a member of the
research team who is not the treating surgeon, to reduce undue influence. The planned
accrual rate, the sentinel staggered enrollment of the 3+3 design, and the Phase 0
readiness gate together give an enrollment period sufficient to treat 18 participants
from approximately 36 screened \cite{daraxwhipple,phase1guidance}.

Sentinel staggering is a binding feature of the first-year conduct: no two
participants in a cohort undergo surgery within the same DLT-evaluation interval
until the supervising surgeon and the DSMB have reviewed the preceding case. After
each participant completes the 28-day DLT-evaluation window and before the next
participant in the cohort undergoes surgery, the accumulating safety data are
reviewed; after each dose-level cohort is complete, the DSMB convenes for a formal
interim safety review and an escalation, expansion, hold, or halt decision. The
prespecified halt rules are quantitative and binding: enrollment and escalation halt
for a mandatory DSMB review if the device-related or procedure-related serious-AE
rate reaches or exceeds one third within a cohort, or if two device-related deaths
occur at any point. Additional pause triggers include any unexpected device-related
serious AE, any breach of a hard cyber-physical safety limit (emergency-stop
latency, positional, or force envelope), and any failure of the pre-procedure
safety-matrix verification. The DSMB may halt, pause, modify, or permit continuation,
and a halt for an immediate hazard is implemented at once \cite{daraxwhipple}.

Oversight is delivered through a three-tier structure with distinct cadences and
authorities. The DSMB operates at cohort boundaries and on triggered review and holds
the authority to halt the study; the Independent Safety Monitor operates per
procedure in real time and holds the authority to stop an individual case; and the
Physical AI Safety Review Committee operates on a cadence of at most 90 days and
reviews the device-and-advisory performance, including the advisory-to-gate
concordance and any model-degradation signal. The reviewing IRB provides ethical
oversight, and the sponsor-investigator is ChemicalQDevice (Kevin Kawchak, CEO),
whose conflict of interest is disclosed under 21 CFR part 54 and mitigated precisely
by the independence of these three tiers, so that no halt, stop, or escalation
decision rests with the financially interested party. The structure is designed so
that the finest-grained safety authority (the per-case stop) sits closest to the
patient and the broadest authority (the study halt) sits with the body that sees the
accumulating cohort data \cite{daraxwhipple,clintrustpai,paiautospons}.

The device operates within a phase-graduated staged-autonomy model. In this Phase 1
study the system runs only at the clinician-controlled (teleoperation or supervised)
level and the supervised semiautonomous level, both under continuous human oversight
with an immediate manual override; full autonomy is prohibited at this phase
consistent with 21 CFR \S 312.21(e). Any future move to higher autonomy (Stage 2
supervised semiautonomous expansion or Stage 3 higher autonomy) is reserved for a
later phase and gated by independent safety review, a USL reassessment, and DSMB
concurrence. The perioperative daraxonrasib pause-and-restart advisory is exercised
in every case: daraxonrasib is paused before surgery and the restart day is set by
an LLM-bound, human-confirmed advisory whose inputs are the postoperative pancreatic
fistula grade by ISGPS (A, B, or C) and the serum daraxonrasib trough relative to a
0.5 ng/mL threshold, with outputs of T+7 days, T+14 days, T+21 days, or hold. A
32-iteration deterministic sweep returned a baseline serum of 0.45 ng/mL
(sub-threshold) in all iterations, a T+7-day restart in 29 of 32 (90.6 percent,
Grade A and sub-threshold), a T+14-day restart in 3 of 32 (9.4 percent, Grade B or
delayed), and no T+21-day restart \cite{daraxwhipple,clintrustpai}.

Safety reporting during the first year follows the 21 CFR \S 312.32 expedited clocks,
augmented by the Physical AI triggers. A fatal or life-threatening suspected adverse
reaction is reported within 7 calendar days and other serious and unexpected
suspected adverse reactions within 15 calendar days. Six additional Physical AI
reporting triggers under \S 312.32(g) apply: a serious Physical AI adverse event on
the 7 or 15-day clock; a system-drug interaction within 15 days; a cybersecurity
incident within 15 days, or within 7 days where patient harm is possible; a sustained
AI model degradation across at least three consecutive procedures or a cumulative
24-hour interval; a sim-to-real divergence exceeding twice tolerance (a trajectory
gap above 4 mm or a force gap above 1.0 N); and a digital-twin discrepancy. For every
Physical AI adverse event an audit trail is preserved from 24 hours before to 72
hours after the event under 21 CFR part 11. Seven Physical AI record types are
retained for a minimum of two years after approval or discontinuation, with HIPAA
Safe Harbor de-identification applied to any disclosed record \cite{cfr312paiuni,congress9510}.

The first-year objectives map onto the endpoint hierarchy of Figure 11: the
co-primary safety, dose-finding, and feasibility readouts are accrued as cohorts
complete, the secondary surgical-quality endpoints are estimated against the Dutch
2025 benchmark as the per-protocol set matures, and the exploratory survival,
advisory-concordance, and sim-to-real measures begin to accumulate through the day-90
pathology review and the 24-month survival follow-up. The realistic first-year
deliverable is completion of the DL1 and DL2 cohorts and initiation of the DL3
cohort, with an interim determination of the MTD and a provisional RP2D if DL3 is
reached and tolerated, together with the device-feasibility and telemetry summaries
that support progression to a later-phase efficacy study \cite{daraxwhipple}.

A representative first-year accrual scenario makes the timeline concrete without
overstating precision. At an accrual of approximately one to two participants per
month, and with the sentinel constraint that each participant's 28-day DLT window
must close (and the supervising surgeon and DSMB must review the case) before the
next participant in the cohort proceeds to surgery, a cohort of three DLT-evaluable
participants occupies on the order of three to four months including the staggering
intervals, and an expanded cohort of six occupies correspondingly longer. Under the
most favorable realization, in which no level requires expansion, the three cohorts
of three traverse DL1, DL2, and DL3 within the first year and yield a provisional
RP2D; under a realization in which one level expands to six, the first year completes
DL1 and DL2 and initiates DL3, which is the realistic deliverable stated above. In
either case the Phase 0 readiness gate (USL at least 7.0, at least 1000 simulations
across at least two frameworks, sim-to-real trajectory below 2 mm and force below
0.5 N) is cleared before the first robotic case, so the first-year clock starts only
after device readiness is documented \cite{phase1guidance,onpremwhippl}.

\subsection{Number of Subjects to be Evaluated}

The planned sample size is up to 18 treated participants, allocated across three
daraxonrasib dose levels (DL1 through DL3) by the 3+3 rule, from approximately 36
screened. This number is not derived from a power calculation, because the study
tests no confirmatory efficacy hypothesis; it is determined jointly by the
dose-finding design and by the early-feasibility precedent for significant-risk
devices, and it is justified by its operating characteristics. The 3+3 rule enrolls
cohorts of three: a level with no DLT in three escalates, a level with one DLT in
three expands to six, and a level with two or more DLTs in up to six is exceeded,
defining the dose below it as the MTD. The maximum exposure at any single dose level
is therefore six participants, and three escalation levels give a design maximum of
three times six, equal to 18 treated participants; many realizations terminate with
fewer. This ceiling is consistent with the FDA early-feasibility convention of
generally fewer than ten participants per device-feasibility cohort while providing
the per-dose exposure the drug-arm 3+3 escalation requires \cite{phase1guidance}.

Rather than power, the design is characterized by its operating characteristics and
by the precision of estimation at the planned sample sizes. The probability that the
3+3 rule escalates past a given dose level is a function of the true DLT probability
at that level: a dose with a true DLT rate near 10 percent is very likely to be
tolerated and escalated, whereas a dose with a true DLT rate of 40 percent or higher
is very likely to be identified as excessively toxic and to define the MTD at the
level below. The exact-interval precision is summarized in Table 4.7. Zero events in
three yields an exact two-sided 95 percent confidence interval of approximately 0 to
71 percent; zero in six, approximately 0 to 46 percent; zero in 18, approximately 0
to 18 percent; and one event in 18, a point estimate of 5.6 percent with an interval
of approximately 0.1 to 27 percent. These intervals make explicit both the value and
the limits of a first-in-human cohort: a clean safety record across the full cohort
of 18 bounds the true event rate below roughly one in five, while a single event
still excludes rates above roughly one quarter \cite{phase1guidance,ichpaistduni}.

\vspace{0.3cm}

\needspace{12\baselineskip}
\noindent\begin{tabularx}{\textwidth}{L{2.9cm} L{3.7cm} Y}
\toprule
\textbf{Quantity} & \textbf{Planned value} & \textbf{Operating characteristic / precision} \\
\midrule
Cohort size & 3, expand to 6 & 3+3 rule: escalate on 0/3; expand on 1/3; exceed dose on $\geq$ 2/6. \\
\addlinespace
Per-dose maximum & 6 & Maximum exposure at any one dose level before the MTD is declared. \\
\addlinespace
Dose levels & 3 (DL1 to DL3) & Three-level escalation; design maximum 3 x 6 = 18 treated. \\
\addlinespace
Total treated & up to 18 (approx. 36 screened) & At or below early-feasibility convention of generally fewer than 10 per device cohort, summed across dose levels. \\
\addlinespace
0 events / 3 & point estimate 0\% & Exact two-sided 95\% CI approximately 0\% to 71\%. \\
\addlinespace
0 events / 6 & point estimate 0\% & Exact two-sided 95\% CI approximately 0\% to 46\%. \\
\addlinespace
0 events / 18 & point estimate 0\% & Exact two-sided 95\% CI approximately 0\% to 18\%. \\
\addlinespace
1 event / 18 & point estimate 5.6\% & Exact two-sided 95\% CI approximately 0.1\% to 27\%. \\
\addlinespace
Significance level & $\alpha$ = 0.05 (two-sided) & Applies to supportive estimation only; no $\alpha$ is spent on dose selection. \\
\bottomrule
\end{tabularx}
\vspace*{-0.3cm}
\figcaption{Table 4.7. Operating characteristics and exact-interval (Clopper-Pearson)
precision at the planned sample sizes. The dose-finding behavior is that of the
standard 3+3 rule; no statistical power is computed because the study tests.}

The relationship between the true DLT rate and the escalation behavior of the 3+3
rule can be made quantitative, and Table 4.8 illustrates it for the first cohort of
three at a level. If the true DLT probability at a level is p, the probability that
none of three participants experiences a DLT (and the level escalates outright) is
the cube of one minus p; for p near 0.10 this is approximately 0.73, for p near 0.20
approximately 0.51, and for p near 0.40 approximately 0.22. The complementary
probability that at least one of three experiences a DLT (and the level either
expands to six or is exceeded) rises correspondingly. The table shows why the 3+3
rule is conservative at low true rates and decisive at high true rates: a level whose
true DLT rate is genuinely near one in ten is escalated most of the time, whereas a
level whose true DLT rate is near two in five triggers expansion or exceedance most
of the time, which is the desired behavior for a first-in-human dose-finding design
\cite{phase1guidance,ichpaistduni}.

\needspace{10\baselineskip}
\noindent\begin{tabularx}{\textwidth}{C{2.8cm} C{3.2cm} C{3.2cm} Y}
\toprule
\textbf{True DLT rate p} & \textbf{P(0 of 3): escalate} & \textbf{P($\geq$1 of 3): expand or exceed} & \textbf{Interpretation} \\
\midrule
0.10 & $\approx$ 0.73 & $\approx$ 0.27 & A genuinely low-toxicity level escalates most of the time. \\
\addlinespace
0.20 & $\approx$ 0.51 & $\approx$ 0.49 & An intermediate level expands or escalates about evenly. \\
\addlinespace
0.30 & $\approx$ 0.34 & $\approx$ 0.66 & A higher-toxicity level usually triggers expansion or exceedance. \\
\addlinespace
0.40 & $\approx$ 0.22 & $\approx$ 0.78 & Toxic level rarely escalates outright; MTD likely set below it. \\
\bottomrule
\end{tabularx}
\vspace*{-0.3cm}
\figcaption{Table 4.8. Escalation behavior of the 3+3 rule for the first cohort of
three as a function of the true dose-limiting-toxicity rate p, computed as the
binomial probability of zero events in three. The rule escalates a genuinely
low-toxicity level most of the time and triggers expansion or exceedance at a toxic
level most of the time.}

The four analysis populations apportion the 18 participants across the analyses
without altering the total exposure: the Safety population comprises all who received
any intervention and is the primary population for safety and primary-endpoint
analyses; the DLT-evaluable population comprises those who completed the 28-day
window or had a DLT within it and drives the escalation and MTD determination; the
Per-Protocol population supports the secondary surgical-quality analyses; and the
modified intention-to-treat population supports the supportive and exploratory
efficacy summaries. Because the cohort is small and intensively monitored,
individual participant data are tabulated in full, so that every aggregate statistic
can be traced to its contributing records \cite{ichpaistduni}.

\subsection{Drug Related Risks}

The foreseeable risks of the investigation arise from three sources, the
investigational drug, the investigational device, and the surgical procedure, and
each is bounded by a specific control rather than left to clinical judgment alone.
The anticipated risks of daraxonrasib are those characteristic of RAS-pathway
inhibition and of the RASolute 302 experience, including gastrointestinal effects
(nausea, diarrhea, and stomatitis), dermatologic effects (rash), hepatic enzyme
elevation, and the class effects of MAPK-pathway suppression; in the perioperative
setting the additional concern is impaired wound healing and anastomotic integrity,
which is the explicit reason for the pause-and-restart advisory that holds
daraxonrasib before surgery and resumes it only on a human-confirmed,
fistula-grade-and-trough-informed schedule (T+7, T+14, T+21 days, or hold). Drug
toxicity is captured through the 28-day DLT window, the dose-blinded SAE attribution,
the exact-interval safety reporting, and the binding cohort-level halt rules
\cite{daraxwhipple,oncotrialllm}. Table 4.9 summarizes the principal risks, their
sources, and their mitigations.

\needspace{14\baselineskip}
\noindent\begin{tabularx}{\textwidth}{L{2.8cm} L{5.4cm} Y}
\toprule
\textbf{Risk source} & \textbf{Foreseeable risk} & \textbf{Mitigation} \\
\midrule
Daraxonrasib (drug) & GI effects (nausea, diarrhea, stomatitis), rash, hepatic enzyme elevation, MAPK-pathway class effects; impaired perioperative wound healing. & 28-day DLT window; conservative 160 mg starting dose below the 300 mg RP2D; pause-and-restart advisory keyed to ISGPS grade and 0.5 ng/mL trough; cohort-level halt rules. \\
\addlinespace
Robotic system (device) & Device malfunction, unsafe trajectory, excessive tip force, vascular injury near the SMV, PV, HA, CA, or SMA. & Force caps $\leq$ 3 N per arm and $\leq$ 18 N cumulative; cross-arm E-stop $\leq$ 3 ms; hardware E-stop $\leq$ 500 ms; 10 kHz heartbeat with 100 microsecond watchdog; vascular no-fly gating; Phase 0 gate (USL $\geq$ 7.0, $\geq$ 1000 sims, sim-to-real $<$ 2 mm and $<$ 0.5 N). \\
\addlinespace
LLM advisory layer & Erroneous or unverified advisory output; over-reliance (therapeutic misconception); model degradation. & Second-opinion advisory only; full autonomy prohibited (\S 312.21(e)); ten-gate verification-before-generation (ACCEPT/BLOCK/ESCALATE); hash-pinned deterministic seed and commit; Physical AI opt-out; advisory-to-gate concordance monitored. \\
\addlinespace
Whipple procedure & Pancreatic fistula, hemorrhage, intra-abdominal infection, sepsis, failure to rescue, unplanned conversion. & Sentinel staggered enrollment with case-by-case DSMB review; three-tier oversight (DSMB, Independent Safety Monitor, Physical AI Safety Review Committee); conversion pathway with continuous human override. \\
\addlinespace
Schedule and population & Delay to effective therapy in a low-survival disease. & Verified single-prompt authoring compresses only the administrative bucket, advancing first dose without altering tumor biology or fixed review clocks (Figure 10). \\
\bottomrule
\end{tabularx}
\vspace*{-0.315cm}
\figcaption{Table 4.9. Principal drug-related, device-related, procedural, and
schedule-related risks with their controls.\\Each risk is bounded by a specific,
prespecified mitigation; the residual risk after these controls is judged.}

The device-related and procedural risks are bounded by the cyber-physical safety
envelope and the three-tier oversight structure. The hard caps (per-arm tip force at
or below 3 N, cumulative cross-arm force at or below 18 N, cross-arm emergency stop
within 3 ms to 0.0 N, and hardware emergency stop within 500 ms), the 10 kHz
heartbeat with the 100 microsecond watchdog, and the five-vessel no-fly gating limit
the magnitude of any single excursion, while the Phase 0 gate (USL at least 7.0, at
least 1000 simulations across at least two frameworks, and a sim-to-real gap below 2
mm and 0.5 N) ensures the system is demonstrably ready before any patient-facing
activity. The LLM advisory risks (an erroneous or unverified advisory and the
therapeutic misconception of over-reliance on the system) are bounded by the
advisory-only role, the prohibition on full autonomy under 21 CFR \S 312.21(e), the
ten-gate verification-before-generation suite, the hash-pinned deterministic seed
and commit, the Physical AI opt-out in the consent, and the monitoring of
advisory-to-gate concordance. Safety reporting follows the 21 CFR \S 312.32 clocks
(fatal or life-threatening suspected adverse reactions within 7 calendar days and
other serious and unexpected reactions within 15 calendar days), augmented by the
six Physical AI reporting triggers under \S 312.32(g) and an audit trail preserved
from 24 hours before to 72 hours after each Physical AI adverse event under 21 CFR
part 11 \cite{cfr312paiuni,congress9510}.

The therapeutic misconception merits a specific control because it is the
characteristic ethical hazard of a study that pairs a novel device with an
advisory model. A participant might over-attribute capability to the robotic system
or to the LLM advisory and consent on the mistaken belief that the technology
guarantees a better outcome than standard care. This is addressed at three points:
in the informed consent under 21 CFR part 50 and ICH E6(R3), which explains that the
device operates only at clinician-controlled and supervised semiautonomous levels
with continuous human oversight and that the LLM is a second opinion rather than a
decision-maker; in the Physical AI opt-out under 21 CFR \S 312.60(f), which lets a
participant decline the Physical AI advisory layer without losing access to the
study's surgical and drug components; and in the staged-autonomy explanation, which
makes clear that full autonomy is prohibited at this phase. The consent is described
to the participant by a research-team member who is not the treating surgeon,
reinforcing that the decision is voluntary and free of the treating surgeon's
influence \cite{cfr050paiuni,clintrustpai}.

The benefit-risk balance is favorable for this population and at this stage. The
population is vulnerable, with a lethal malignancy whose five-year survival is below
13 percent and for whom the resection margin is the strongest modifiable determinant
of survival; the intervention pairs a genotype-matched systemic agent with a
margin-negative curative-intent resection delivered within hard cyber-physical limits
under continuous human oversight. The schedule benefit reaches the patient through
the cascade of Figure 10: faster validated authoring advances IND and IRB submission
and site activation, the elapse of the 30-day clock and the first dose, and the 3+3
cohort turnover toward the RP2D, shortening the white space between phases and, for
this low-survival population, extending progression-free and overall survival,
without changing tumor biology or any fixed review clock. For patients whose disease
will not pause, the alternative of slower sequential authoring carries its own,
often larger, cost in delayed treatment. With each foreseeable risk bounded by a
prespecified control, with the conservative 160 mg starting dose, the sentinel
staggered enrollment, the binding halt rules, and the three-tier oversight, the
probable benefit of the investigation exceeds its probable risk, and the General
Investigational Plan supports the proposed first-year conduct \cite{pdac060s2030,daraxwhipple,clintrustpai}.

\subsection{References}

The clinical, statistical, and acceleration content of this plan is grounded in the
trial protocol and the supporting evidence base. The Phase 1 first-in-human guidance
and the early-feasibility framework inform the single-arm dose-finding design and the
operating characteristics \cite{phase1guidance}. The daraxonrasib robotic Whipple
program and the on-premises eight-arm system specifications supply the drug, device,
and advisory facts \cite{daraxwhipple,onpremwhippl,oncotrialllm}. The Physical AI
adaptations of 21 CFR Part 312 and Part 50 and the ICH guidance govern the safety
reporting, consent, and statistical conventions \cite{cfr312paiuni,cfr050paiuni,ichpaistduni}.
H.R. 9510 and the clinical-trust framework establish the verification-before-generation
discipline \cite{congress9510,clintrustpai}. The PDAC epidemiology, the acceleration
mechanism, and the supporting computational evidence are drawn from the population
and document-acceleration sources \cite{pdac060s2030,paipredict4x,natlpaiplatf,paisitedocpk,paiautospons,llmcomp16fda,fdadigtwinpc,qspmetpancre,chatgpt100kp,pdacdigtwinp}.
The full bibliography is consolidated in the references and back-matter section of
this IND.
